% Deadline contracts: sample construction, measurement and tests.
% All numbers come from the saved study outputs under data/research/pm_hazard/ (summary.json,
% scan_plan.parquet) and the study report data/reports/research-pm-hazard-20261109T222119Z/report.md.

A deadline contract pays one dollar if a specified occurrence happens after issuance and before a
stated date $T$; the exchange may close it as soon as the occurrence is confirmed. Under the model of
\cref{sec:model} its price is one minus a survival probability. We measure three objects: the implied
average hazard $\hbarimp_t$, its term structure across deadlines, and the no-event decay (theta) of
the price.

\paragraph{Sample and classification.}
Deadline contracts live for weeks or months, so we extend the settled-market table of \cref{sec:data}
backwards series by series through the public historical endpoint. The extension covers 237 series
with a deadline contract in the table and 1{,}450 further non-financial one-off, custom or periodic
series whose title contains a deadline-type word (``before'', ``out'', ``release'', ``visit'',
``eliminated'' and similar). Series are selected on titles, never on outcomes.
% source: data/research/pm_hazard/scan_plan.parquet (tier 1: 237, tier 2: 1450)
Contracts are classified from their rules text alone. A candidate's rules contain a deadline phrase
(``before $D$'', ``by $D$'', ``through $D$'', ``between $D_1$ and $D_2$'', optionally with a clock
time), and candidates fall into four hazard types (\cref{haz-tab-class}):
\begin{itemize}
\item \emph{events}: single occurrences, such as a release, resignation, visit or launch;
\item \emph{mentions}: weekly or monthly windows in which a named speaker must say a word;
\item \emph{counts}: first-passage thresholds (``at least $N$ \dots{} before $D$'');
\item \emph{eliminations}: reality-show eliminations, resolved at scheduled episodes.
\end{itemize}
We exclude competing legs, which resolve NO early when another candidate succeeds; quantities
measured at $T$; scheduled votes; and scheduled shows and elections. A contract is in scope when its
parsed deadline lies in $(1\text{ January}, 2\text{ November }2026]$, it settled YES or NO, and it
was open after 1 January 2026. The upper bound is the key look-ahead guard: a contract whose deadline
lies beyond the data end would appear only if it had already resolved YES. The sample has 5{,}872
contracts, 5{,}761 of them with hourly quotes.
% source: data/research/pm_hazard/summary.json:universe (in_scope_markets 5872, with_candles 5761)

\begin{table}[t]
\centering
\begin{threeparttable}
\caption{Classification of deadline contracts}\label{haz-tab-class}
\small
% source: data/reports/research-pm-hazard-20261109T222119Z/report.md (classification table)
\setlength{\tabcolsep}{4pt}
\begin{tabular}{@{}lrrrrrp{4.8cm}@{}}
\toprule
Type & Cand. & Series & In scope & Clusters & YES rate & Example \\
\midrule
Event        & 3{,}505 & 700 & 2{,}180 & 548 & 0.306 & Google releases next Gemini Pro before 16 Oct \\
Mention      & 2{,}567 &   8 & 1{,}940 &  71 & 0.588 & Trump says ``Epstein'' before 2 Nov \\
Elimination  & 1{,}485 &   8 & 1{,}449 &  88 & 0.109 & Contestant eliminated from DWTS before 28 Oct \\
Count        &    504  &  31 &    303  &  52 & 0.399 & Trump endorses $\ge 5$ people this week \\
\midrule
\multicolumn{7}{@{}l}{\emph{Excluded}} \\
Competing legs     & 2{,}401 & 64 & & & & Player's next team before 1 Dec: Washington \\
Measured at $T$    & 1{,}014 & 25 & & & & Weekly token total above a level \\
Scheduled vote     &    250  & 10 & & & & Senator votes for a cloture motion \\
Scheduled show or election & 36 & 7 & & & & Artist performs a song at a named show \\
\bottomrule
\end{tabular}
\begin{tablenotes}\footnotesize
\item Candidates (Cand.) are non-financial single-leg contracts whose rules contain a deadline
phrase. In scope: parsed deadline in (1 January, 2 November 2026], settled YES or NO, open after
1 January 2026. Clusters: distinct exchange event tickers among in-scope contracts (the clustering
unit). YES rate: share of in-scope contracts that resolved YES. Examples paraphrase contract titles.
\end{tablenotes}
\end{threeparttable}
\end{table}

\paragraph{Deadlines, quotes and look-ahead guards.}
The deadline is parsed from the rules: ``before $D$'' ends at 00:00 US Eastern on $D$; ``by'',
``through'' and ``between \dots{} and $D$'' include $D$; an explicit clock time overrides. The window
start $s$ is an explicit ``after $D$''/``from $D$'' date or, failing that, issuance. Close and
settlement times are never used, since they reveal an early YES. Of contracts resolving NO, 97.6\%
closed within 26 hours of the parsed $T$ (median offset $-0.02$ hours). Contracts resolving YES
closed more than six hours early in 90.1\% of events and 81.8\% of mentions, but only 31.6\% of
eliminations.
% source: summary.json:deadline_validation
Decisions are taken at $t=T-L$ for leads $L$ of 1, 3, 6, 12, 24, 48, 72, 120, 168, 240, 336, 504,
720, 1080 and 1440 hours. The quote is the bid and ask of the last hourly candle that ended strictly before $t$, at most
two hours old. An observation is \emph{priced} when the spread is at most 10\cents{} (fixed in
advance). The \estim{} (before 1 October 2026) holds contracts with $T$ before 1 October 2026. The
holdout (October 2026) holds decisions taken in October with $T\le 2$ November. Straddling contracts
enter pooled tables only. The 26{,}290 priced observations cover 5{,}070 contracts in 709 event
clusters: 22{,}736 in the \estim{} and 2{,}076 in the holdout.
% source: summary.json:universe
Measures that condition on ``no event yet'' use an upper bound on the event time, $\tau_E$: the
earlier of the YES close and the end of the first hourly candle of the final run of mids of at least
97\cents. $\tau_E$ never enters calibration or trading statistics. Conditioned measures are
restricted to mids below 0.9, so that a confidently priced contract is not flagged as resolved before
its event.

\paragraph{Implied hazards and strips.}
$\hbarimp_t=-\ln(1-p_t)/(T-\max(t,s))$, with $p_t$ the mid. A strip is a set of event or mention
contracts with identical rules apart from dates. Legs are compared only while both are open and
their windows have started, and only if the later deadline does not start later than the earlier one.
Where a window starts at issuance, this assumes that a later deadline issued no later than an earlier
one covers its window; without such nesting, monotonicity in $T$ is not an arbitrage bound. From one
daily snapshot per strip (16:00 UTC, spread at most 10\cents) we bootstrap forward hazards
(\cref{sec:model}). The shape is the ratio of the last forward hazard to the spot hazard to the first
deadline: rising above 1.25, front-loaded below 0.8, otherwise flat, and inverted if any forward is
negative.

\paragraph{Calibration ratio.}
In a cell of observations with mids $p_i$ and outcomes $y_i$, the hazard multiplier $k$ solves
\begin{equation}
\textstyle\sum_i\bigl[1-(1-p_i)^{k}\bigr]=\sum_i y_i .\label{haz-eq-k}
\end{equation}
Because $1-(1-p)^k=1-\exp(-k\,\hbarimp\tau)$, $k$ scales every implied cumulative hazard and, under a
flat hazard, is the ratio of realised to implied hazard; $k<1$ means implied hazards are too high.
Intervals are the 5th and 95th percentiles over 200 bootstrap resamples of event clusters
\citep{camerongelbachmiller2008}. We also report the mean of $y-p$ in cents with event-clustered
$t$-statistics (\cref{sec:method}).

\paragraph{Theta and event timing.}
Each priced observation is followed to the next grid point $t_2$, which is set by the grid whether or
not the contract is still open ($t_2=T$ at the 1-hour lead); the step contains a jump if
$\tau_E\le t_2$. We compare the realised no-jump change with the flat-hazard theta
$(1-p_t)(1-e^{\hbarimp_t\Delta})$, the jump count with $\sum(1-e^{-\hbarimp_t\Delta})$, and report
the martingale drift $\E[p_{t_2}-p_t]$, with $p_{t_2}=1$ after a jump. Decay paths compare $p(L)/p(A)$
for contracts with no event yet at an anchor $A=168$ hours against the flat-hazard path
$[1-(1-p_A)^{\tau_L/\tau_A}]/p_A$. For timing, each YES contract is anchored at its earliest priced
pre-event quote $p_0$ at $t_0$, with $s_0=\max(t_0,s)$ and elapsed window share
$u=(\tau_E-s_0)/(T-s_0)$. Under the flat hazard implied by $p_0$, $F(u)=[1-(1-p_0)^u]/p_0$ is uniform
\citep{diebold1998}; a mean above 0.5 means events come later than implied.

\paragraph{Jumps, diffusion and the Cox intensity.}
The quadratic variation of the hourly mid (spread at most 10\cents) is measured from the first quote in
the last 62 days to resolution and split four ways:
\begin{itemize}
\item diffusion: hourly changes below 20\cents{} before resolution;
\item pre-resolution jumps: hourly changes of 20\cents{} or more before resolution;
\item the YES resolution jump $(1-m)^2$, from the last mid at least two hours before $\tau_E$;
\item the NO deadline step $m^2$.
\end{itemize}
This follows the separation of continuous and jump variation in \citet{barndorffnielsen2004} and
\citet{andersenbollerslevdiebold2007}. For a bounded martingale settling at 0 or 1 the expected total
is $p_0(1-p_0)$; bid-ask bounce inflates the realised value \citep{roll1984,zhang2005}. On no-jump
steps with mids in $(0.02,0.98)$ we fit a CIR-type intensity \citep{cox1985cir,duffiepansingleton2000}
to $\lambda=\hbarimp$ by least squares without intercept,
\begin{equation}
\frac{\Delta\lambda}{\sqrt{\lambda}}=\kappa\theta\,\frac{\Delta t}{\sqrt{\lambda}}-\kappa\sqrt{\lambda}\,\Delta t+\sigma\sqrt{\Delta t}\,\varepsilon ,\label{haz-eq-cir}
\end{equation}
separately for leads above and below three days. Because $\hbarimp$ is an average hazard to $T$, a
deadline-shaped intensity loads on $\kappa$ and $\theta$, so the fit is descriptive.

\paragraph{Static arbitrage and economic significance.}
For adjacent nested legs quoted in the same hourly candle, we count the pair-hours in which the
earlier deadline's mid exceeds the later one's, and those in which its bid exceeds the later ask,
before and after the taker fee on both legs. Buying YES on the later deadline and NO on the earlier
one pays at least one dollar in every state; displayed sizes are not observed. The trading test is
recorded briefly. A multiplier table by type, lead and price is fitted on the \estim, and the rule
buys NO at the NO ask when the model's edge net of fees exceeds a threshold chosen on the \estim.
It is compared with a price-only calibration model and with a control that buys NO on every quote
below a dev-chosen price.
